\begin{thebibliography}{50}
\bibitem{_Ref213949315}
Zhang W, Zhao X, Chen Y, et al. DeepUWF: an automated ultra-wide-field fundus screening system via deep learning. IEEE Journal of Biomedical and Health Informatics 2020;25(8): 2988-2996.
\bibitem{_RefRahman2022MLFundus}
Mujeeb Rahman K K, Nasor M, Imran A. Automatic screening of diabetic retinopathy using fundus images and machine learning algorithms. Diagnostics 2022;12(9): 2262.
\bibitem{_RefLim2024VLVOphthalmology}
Lim G, Elangovan K, Jin L. Vision language models in ophthalmology. Current Opinion in Ophthalmology 2024;35(6): 487-493.
\bibitem{_RefUrooj2025LLMMedicalImage}
Urooj B, Fayaz M, Ali S, et al. Large language models in medical image analysis: a systematic survey and future directions. Bioengineering 2025;12(8): 818.
\bibitem{_Ref213871637}
Kumar V, Singh R S, Rambabu M, et al. Deep learning for hyperspectral image classification: A survey. Computer Science Review 2024;53: 100658.
\bibitem{_Ref213948625}
Mo Y, Wu Y, Yang X, et al. Review the state-of-the-art technologies of semantic segmentation based on deep learning. Neurocomputing 2022;493: 626-646.
\bibitem{_Ref213948898}
Qiu D, Cheng Y, Wang X. Medical image super-resolution reconstruction algorithms based on deep learning: A survey. Computer Methods and Programs in Biomedicine 2023;238: 107590.
\bibitem{_Ref213949262}
Hasan S M M, Uddin M P, Al Mamun M, et al. A machine learning framework for early-stage detection of autism spectrum disorders. IEEE Access 2022;11: 15038-15057.
\bibitem{_Ref213949362}
Al-Fahdawi S, Al-Waisy A S, Zeebaree D Q, et al. Fundus-deepnet: Multi-label deep learning classification system for enhanced detection of multiple ocular diseases through data fusion of fundus images. Information Fusion 2024;102: 102059.
\bibitem{_Ref213949382}
Chen Z M, Wei X S, Wang P, et al. Multi-label image recognition with graph convolutional networks. In: Proceedings of the IEEE/CVF conference on computer vision and pattern recognition; 2019. p. 5177-5186.
\bibitem{_Ref213949348}
Liu S, Yin S, Qu L, et al. Reducing domain gap in frequency and spatial domain for cross-modality domain adaptation on medical image segmentation. In: Proceedings of the AAAI conference on artificial intelligence; 2023;37(2): 1719-1727.
\bibitem{_Ref213949354}
Yang C, Guo X, Chen Z, et al. Source free domain adaptation for medical image segmentation with fourier style mining. Medical Image Analysis 2022;79: 102457.
\bibitem{_Ref213949393}
Bertino E, Ooi B C, Yang Y, et al. Privacy and ownership preserving of outsourced medical data. In: 21st International Conference on Data Engineering (ICDE\textquotesingle05). IEEE; 2005. p. 521-532.
\bibitem{_Ref213949872}
Wang Y, Xie Y, Fan L, et al. STMG: Swin transformer for multi-label image recognition with graph convolution network. Neural Computing and Applications 2022;34(12): 10051-10063.
\bibitem{_Ref213949896}
Zhang M L, Li Y K, Liu X Y, et al. Binary relevance for multi-label learning: an overview. Frontiers of Computer Science 2018;12(2): 191-202.
\bibitem{_Ref213949927}
Ridnik T, Ben-Baruch E, Zamir N, et al. Asymmetric loss for multi-label classification. In: Proceedings of the IEEE/CVF international conference on computer vision; 2021. p. 82-91.
\bibitem{_Ref213950246}
Liu S, Zhang L, Yang X, et al. Query2label: A simple transformer way to multi-label classification. arXiv preprint arXiv:2107.10834. 2021.
\bibitem{_Ref213950282}
Everingham M, Van Gool L, Williams C K I, et al. The pascal visual object classes (voc) challenge. International journal of computer vision 2010;88(2): 303-338.
\bibitem{_Ref213950349}
Cheng Y, Ma M, Li X, et al. Multi-label classification of fundus images based on graph convolutional network. BMC Medical Informatics and Decision Making 2021;21(Suppl 2): 82.
\bibitem{_Ref213950338}
He J, Li C, Ye J, et al. Multi-label ocular disease classification with a dense correlation deep neural network. Biomedical Signal Processing and Control 2021;63: 102167.
\bibitem{_Ref213950343}
Bhati A, Gour N, Khanna P, et al. Discriminative kernel convolution network for multi-label ophthalmic disease detection on imbalanced fundus image dataset. Computers in Biology and Medicine 2023;153: 106519.
\bibitem{_Ref213949439}
Guan H, Liu M. Domain adaptation for medical image analysis: a survey. IEEE Transactions on Biomedical Engineering 2021;69(3): 1173-1185.
\bibitem{_Ref213949448}
Tzeng E, Hoffman J, Zhang N, et al. Deep domain confusion: Maximizing for domain invariance. arXiv preprint arXiv:1412.3474. 2014.
\bibitem{_Ref213949453}
Sun B, Feng J, Saenko K. Return of frustratingly easy domain adaptation. In: \allowbreak Proceedings of the AAAI conference on artificial intelligence; 2016;30(1).
\bibitem{_Ref213949479}
Ganin Y, Ustinova E, Ajakan H, et al. Domain-adversarial training of neural networks. Journal of machine learning research 2016;17(59): 1-35.
\bibitem{_Ref213949490}
Zhu J Y, Park T, Isola P, et al. Unpaired image-to-image translation using cycle-consistent adversarial networks. In: Proceedings of the IEEE international conference on computer vision; 2017. p. 2223-2232.
\bibitem{_Ref213949530}
Wang R, Wu Z, Weng Z, et al. Cross-domain contrastive learning for unsupervised domain adaptation. IEEE Transactions on Multimedia 2022;25: 1665-1673.
\bibitem{_Ref213949554}
Zhou K, Liu Z, Qiao Y, et al. Domain generalization: A survey. IEEE transactions on pattern analysis and machine intelligence 2022;45(4): 4396-4415.
\bibitem{_Ref213949562}
Zhou K, Yang Y, Qiao Y, et al. Domain generalization with mixstyle. arXiv preprint arXiv:2104.02008. 2021.
\bibitem{_Ref213949607}
Gulrajani I, Lopez-Paz D. In search of lost domain generalization. arXiv preprint arXiv:2007.01434. 2020.
\bibitem{_Ref213949788}
Xu Q, Zhang R, Zhang Y, et al. A fourier-based framework for domain generalization. In: Proceedings of the IEEE/CVF conference on computer vision and pattern recognition; 2021. p. 14383-14392.
\bibitem{_Ref213949796}
Huang Z, Wang H, Xing E P, et al. Self-challenging improves cross-domain generalization. In: European conference on computer vision. Cham: Springer International Publishing; 2020. p. 124-140.
\bibitem{_Ref213949803}
Khoee A G, Yu Y, Feldt R. Domain generalization through meta-learning: a survey. Artificial Intelligence Review 2024;57(10): 285.
\bibitem{_Ref213949863}
Ouyang C, Chen C, Li S, et al. Causality-inspired single-source domain generalization for medical image segmentation. IEEE Transactions on Medical Imaging 2022;42(4): 1095-1106.
\bibitem{_Ref213950358}
Yadav N, Pandey S, Gupta A, et al. Data privacy in healthcare: in the era of artificial intelligence. Indian Dermatology Online Journal 2023;14(6): 788-792.
\bibitem{_Ref213950376}
Aledhari M, Razzak R, Parizi R M, et al. Federated learning: A survey on enabling technologies, protocols, and applications. IEEE Access 2020;8: 140699-140725.
\bibitem{_Ref213950391}
Li T, Sahu A K, Zaheer M, et al. Federated optimization in heterogeneous networks. Proceedings of Machine learning and systems 2020;2: 429-450.
\bibitem{_Ref213950427}
Martins P, Sousa L, Mariano A. A survey on fully homomorphic encryption: An engineering perspective. ACM Computing Surveys (CSUR) 2017;50(6): 1-33.
\bibitem{_Ref213950433}
Dutil F, See A, Di Jorio L, et al. Application of homomorphic encryption in medical imaging. arXiv preprint arXiv:2110.07768. 2021.
\bibitem{_Ref213950438}
Zhao C, Zhao S, Zhao M, et al. Secure multi-party computation: theory, practice and applications. Information Sciences 2019;476: 357-372.
\bibitem{_Ref213950443}
Bu Z, Dong J, Long Q, et al. Deep learning with gaussian differential privacy. Harvard data science review 2020;2020(23): 10.1162/99608f92.cfc5dd25.
\bibitem{_Ref213950449}
Dong J, Roth A, Su W J. Gaussian differential privacy. Journal of the Royal Statistical Society Series B: Statistical Methodology 2022;84(1): 3-37.
\bibitem{_Ref213950458}
Bagdasaryan E, Poursaeed O, Shmatikov V. Differential privacy has disparate impact on model accuracy. Advances in neural information processing systems 2019;32.
\bibitem{_RefResNet2016}
He K, Zhang X, Ren S, et al. Deep residual learning for image recognition. In: Proceedings of the IEEE conference on computer vision and pattern recognition; 2016. p. 770-778.
\bibitem{_Ref214903493}
Dwork C, McSherry F, Nissim K, et al. Calibrating noise to sensitivity in private data analysis. In: Theory of cryptography conference. Berlin, Heidelberg: Springer Berlin Heidelberg, 2006: 265-284.
\bibitem{_Ref213949402}
Li N, Li T, Hu C, et al. A benchmark of ocular disease intelligent recognition: One shot for multi-disease detection. In: International symposium on benchmarking, measuring and optimization. Cham: Springer International Publishing; 2020. p. 177-193.
\bibitem{_Ref213949409}
Pachade S, Porwal P, Thulkar D, et al. Retinal fundus multi-disease image dataset (rfmid): A dataset for multi-disease detection research. Data 2021;6(2): 14.
\bibitem{_RefEDID2024}
Sharmin S, Rashid M R, Khatun T, et al. A dataset of color fundus images for the detection and classification of eye diseases. Data in Brief 2024;57: 110979.
\bibitem{_Ref214903299}
Gour N, Khanna P. Multi-class multi-label ophthalmological disease detection using transfer learning based convolutional neural network. Biomedical signal processing and control 2021;66: 102329.
\bibitem{_Ref214903327}
Ou X, Gao L, Quan X, et al. BFENet: A two-stream interaction CNN method for multi-label ophthalmic diseases classification with bilateral fundus images. Computer Methods and Programs in Biomedicine 2022;219: 106739.
\end{thebibliography}
